\clearpage\fancyhead[L]{\small PALLAS.}\hypertarget{native-336-lower}{}\label{native-336-lower}
\begin{center}\textbf{PALLAS.}\par\vspace{4pt}\hrule\vspace{4pt}\fontsize{10}{12}\selectfont Astronomisches Jahrbuch für 1814. Seite 246\,..\,249. Berlin 1811.\par\vspace{4pt}\hrule\vspace{4pt}\end{center}

\begin{flushright}Göttingen, den 10. Februar 1811.\end{flushright}
\par\indent Am 14. Januar dieses Jahres habe ich in unseren gelehrten Anzeigen einen kurzen Bericht über die Störungen der Pallas durch Jupiter veröffentlichen lassen. Die \emph{allgemeine} Theorie dieser Störungen steht freilich noch aus. Sie wird noch viel mehr Arbeit erfordern, wenn sie alles leisten soll, was man sich wünschen kann. Ich habe damit schon angefangen. Es eilt aber noch nicht, da sich die Art, wie ich die Störungen in meiner jüngsten Arbeit behandelt habe, noch viele Jahre in gleicher Weise fortführen lässt. Sobald die Opposition dieses Jahres beobachtet worden ist, werde ich diese Fortsetzung in Angriff nehmen und dann vielleicht auch Juno und Vesta ähnlich behandeln. Ich bitte Sie, die bevorstehende $\text{\AstroSymbols ☍}$ der Pallas sorgfältig zu beobachten. Zur Erleichterung gebe ich hier die verbesserte Ephemeride an, denn die Übereinstimmung meiner neuen Ergebnisse hat sich vollständig bestätigt.
\Needspace{14\baselineskip}
\begin{center}\fontsize{10}{12}\selectfont\setlength{\tabcolsep}{4pt}\renewcommand{\arraystretch}{1.15}\begin{tabular}{lll}
Febr. 11 & AR. $148^\circ\,3'$ & Dekl. $15^\circ\,43'$ S.\\
\phantom{Febr. }15 & $147\;\;20$ & $14\;\;16$\\
\phantom{Febr. }19 & $146\;\;37$ & $12\;\;42$\\
\phantom{Febr. }23 & $145\;\;55$ & $11\;\;3$\\
\phantom{Febr. }27 & $145\;\;17$ & $9\;\;19$\\
März \phantom{1}3 & AR. $144^\circ\,43'$ & Dekl. $7^\circ\,32'$ S.\\
\phantom{März }7 & $144\;\;13$ & $5\;\;44$\\
11 & $143\;\;49$ & $3\;\;57$\\
15 & $143\;\;31$ & $2\;\;11$\\
19 & $143\;\;20$ & $0\;\;29$\\
\end{tabular}\end{center}

\SourceLeaf{337}\Needspace{12\baselineskip}\par\indent Professor \textsc{Harding} hat in Mannheim die folgenden Beobachtungen der Pallas gemacht.
\begin{center}\fontsize{10}{12}\selectfont\setlength{\tabcolsep}{4pt}\renewcommand{\arraystretch}{1.15}
\begin{tabular}{l|c|c|c}
1811 Januar 20 & $14^{\mathrm u}\,6^{\mathrm m}\,38^{\mathrm s}6$ & $151^\circ\,10'\,36''2$ & $-21^\circ\,7'\,39''7$ \\
\phantom{1811 Januar }21 & $14\;\;2\;\;18.1$ & $151\;\;4\;\;53.5$ & $-20\;\;57\;\;27.7$ \\
\phantom{1811 Januar }22 & $13\;\;57\;\;58.4$ & & $-20\;\;47\;\;51.7$ \\
\phantom{1811 Januar }23 & $13\;\;53\;\;38.6$ & $150\;\;52\;\;48.2$ & $-20\;\;37\;\;58.8$ \\
\phantom{1811 Januar }24 & $13\;\;49\;\;16.4$ & $150\;\;46\;\;0.6$ & $-20\;\;27\;\;9.7$ \\
\phantom{1811 Januar }25 & $13\;\;44\;\;53.0$ & $150\;\;39\;\;9.9$ & $-20\;\;16\;\;9.4$ \\
\end{tabular}
\end{center}


\par\indent Der Planet ist jetzt ziemlich hell. Ceres, die sehr nahe an der in ihrer Ephemeride angegebenen Position steht, schien mir aber noch etwas heller zu sein.
\Needspace{18\baselineskip}\begin{flushright}Göttingen, den 29. August 1811.\end{flushright}
\par\indent Verehrtester Freund, ich freue mich, Ihnen anbei die Ephemeriden für die nächste Erscheinung von \emph{Pallas}, \emph{Juno} und \emph{Vesta} zu übersenden. Drei geschickte Schüler von mir haben sich die Rechnung geteilt. Näheres zu Pallas finden Sie in Nr. 8 und 130 der Göttingischen Gelehrten Anzeigen [S. 327 d. B]. Herr \textsc{Nicolai} hat unter meiner Aufsicht die Berechnung der Störung der Pallas durch Jupiter jetzt um ein weiteres Jahr fortgeführt und berechnet nun die Ephemeride für 1813. Näheres zu Juno steht in Nr. 92 unserer gelehrten Anzeigen [Seite 325 d. B.] Die Elemente der Vesta sind noch vorläufig und werden eine kleine Verbesserung brauchen. Inzwischen teile ich Ihnen hier die Beobachtungen mit, die ich an unserer Sternwarte mit dem Mauerquadranten gemacht habe.
\Needspace{9\baselineskip}\begin{center}\fontsize{10}{12}\selectfont\setlength{\tabcolsep}{4pt}\renewcommand{\arraystretch}{1.15}
\begin{tabular}{l|c|c|c}
1811 Mai 20 & $12^{\mathrm u}\,27^{\mathrm m}\,45''1$ M.\,Z. & $244^\circ\,42'\,0''9$ & $-12^\circ\,28'\,10''1$ \\
\phantom{1811 Mai }24 & $12\;\;8\;\;3.4$ & $243\;\;42\;\;26.1$ & $-12\;\;30\;\;1.8$ \\
\phantom{1811 Mai }25 & $12\;\;3\;\;7.6$ & $243\;\;27\;\;25.3$ & $-12\;\;30\;\;47.8$ \\
\phantom{1811 Mai }29 & $11\;\;43\;\;22.9$ & $242\;\;27\;\;0.7$ & $-12\;\;35$ \\
\phantom{1811 Mai }30 & $11\;\;38\;\;27.3$ & $242\;\;12\;\;3.0$ & $-12\;\;35\;\;52.8$ \\
\end{tabular}
\end{center}


\par\indent Aus diesen Beobachtungen, die ich alle für gut halte, habe ich die Opposition wie folgt abgeleitet. 1811 Mai 15. $12^{\mathrm u}\,44^{\mathrm m}\,30^{\mathrm s}$ mittlere Zeit in Göttingen, $243^\circ\,48'\,43''9$ wahre Länge, $8^\circ\,33'\,59''5$ nördliche geozentrische Breite. Zu den Änderungen, die Herr Dr. \textsc{Triesnecker} dem Jahrbuch von 1813 zufolge an meinen Ergebnissen für einige Ceresoppositionen vornehmen zu müssen glaubte, bemerke ich Folgendes.

\par\indent 1) In der Berechnung der Opposition für 1807 aus den Mailänder Beobachtungen habe ich, soweit ich sie bisher überprüft habe, keine Fehler gefunden. Ich muss daher bei meinem Ergebnis vom 3. Mai bleiben, $8^{\mathrm u}\,42^{\mathrm m}\,28^{\mathrm s}$ Pariser Zeit und $222^\circ\,14'\,9''1$ Länge. Ich will aber die ganze Rechnung wiederholen, sobald ich dazu komme, die Elemente nach den beiden Oppositionen von 1809 und 1811 (die noch nicht benutzt worden sind) zu berechnen.

\par\indent 2) Die Vermutung von Herrn Dr. \textsc{Triesnecker}, bei der Opposition von 1808 die geozentrische Breite von $12^\circ\,43'\,44''8$ in $12^\circ\,45'\,44''8$ zu ändern, ist völlig unzulässig. Ebenso muss die Länge der Knotenlinie unbedingt $80^\circ\,53'\,23''$ bleiben und darf keineswegs um $5'$ größer angesetzt werden. Herr Tr. kam nur deshalb zu seiner Vermutung, weil die Elemente mit meinen Störungstafeln sonst weniger gut mit der Breite übereinstimmen. Ich wundere mich aber sehr, dass er meinen ausdrücklichen Hinweis übersehen hat. Bei der Berechnung der Breiten habe ich anstelle meiner alten Tafeln die vollständigen Störungsgleichungen verwendet, die ich später entwickelt habe. Sie sind bisher noch nicht gedruckt worden. Deshalb teile ich sie Ihnen hier so mit, wie sie 1805 berechnet worden sind.
\SourceLeaf{338}\Needspace{24\baselineskip}\begin{center}\emph{Störungen der Breite der Ceres durch Jupiter.}

($\JupiterSymbol$ mittlere Länge des $\JupiterSymbol$, $\CeresSymbol$ mittlere Länge der $\CeresSymbol$)
\end{center}
\[
\begin{array}{l}
-9''65\,\sin(\JupiterSymbol-74^\circ\,13')\\
-14.17\,\sin(\CeresSymbol-2\JupiterSymbol+78^\circ\,35')\\
+27.47\,\sin(2\CeresSymbol-3\JupiterSymbol+78\;35)\\
-4.62\,\sin(3\CeresSymbol-4\JupiterSymbol+78\;35)\\
+5.51\,\sin(2\CeresSymbol-\JupiterSymbol-78\;35)\\
+0.99\,\sin(3\CeresSymbol-2\JupiterSymbol-78\;35)\\
+1.91
\end{array}
\]
\[
\begin{array}{l}
+8.61\,\sin(\CeresSymbol-\JupiterSymbol+95\;26)\\
-10.31\,\sin(2\CeresSymbol-2\JupiterSymbol+94\;10)\\
-6''16\,\sin(3\CeresSymbol-3\JupiterSymbol+112^\circ\,26')\\
-13.69\,\sin(2\JupiterSymbol-57^\circ\,21')\\
-6.77\,\sin(3\JupiterSymbol-\CeresSymbol-61\;25)\\
+4.30\,\sin(2\CeresSymbol-4\JupiterSymbol+72\;24)\\
-22.51\,\sin(3\CeresSymbol-5\JupiterSymbol+70\;27)\\
+1.65\,\sin(3\JupiterSymbol-56^\circ\,9')
\end{array}
\]
\par\indent Die ersten sechs Gleichungen sind die alten in verbesserter Form. Die übrigen neun, die eine Minute betragen können, sind ganz neu hinzugekommen. Daher überrascht es nicht, dass sich die Lage des Knotens bei Berücksichtigung aller Gleichungen um $5'$ änderte. — Diese 15 Breitengleichungen sind noch zu ertragen, aber bei Pallas habe ich über 40, die über $1''$ hinausgehen.\par
\indent Den Kometen haben wir am 22. August zum ersten Mal wieder gesehen. Erst heute früh hat Professor \textsc{Harding} aber einige seiner Abstände von Fixsternen gemessen. Sie sind noch nicht reduziert worden.
